\documentclass[11pt]{article}
\usepackage[margin=1in]{geometry}
\usepackage{amsmath,amssymb,amsthm,mathtools}
\usepackage{graphicx}
\usepackage{microtype}
\usepackage[hidelinks]{hyperref}
\usepackage{enumitem}

\newtheorem{theorem}{Theorem}
\newtheorem{lemma}{Lemma}
\newcommand{\T}{\mathbb T}
\newcommand{\cF}{\mathfrak F}
\newcommand{\cQ}{\mathcal Q}
\newcommand{\cP}{\mathcal P}

\title{A Hopf-Bundle Counterexample to Hatori's\\
Separation Condition for Commutative JB$^*$-Triples}
\author{Automated mathematics run \texttt{fa\_banach\_001}, lane 1}
\date{August 9, 2026}

\begin{document}
\maketitle

\begin{center}
\fbox{\parbox{0.9\textwidth}{\textbf{Status.}
Candidate counterexample, likely valid, pending expert review.  A compact
principal circle-bundle triple $C^{\T}(L)$ satisfies Hatori's condition
$(*)$ exactly when the bundle has a continuous global section.  The Hopf
bundle therefore gives a commutative JB$^*$-triple which fails $(*)$.}}
\end{center}

\section{Source problem}

Hatori~\cite{Hatori} associates to a complex Banach space $E$ a set
$\cQ_E$ of extreme dual-ball functionals which expose the maximal convex
subsets of the unit sphere.  A set of representatives $\cP\subset\cQ_E$
contains one functional from every equivalence class up to unimodular scalar.
His condition $(*)$ says that some such $\cP$ has the following property:
for every $p\in\cP$, every $\varepsilon>0$, and every relatively weak$^*$
closed $F\subset\cP$ with $p\notin F$, there is $a\in S(E)$ such that
\begin{equation}\label{eq:star}
 p(a)=1,\qquad |q(a)|\le\varepsilon\quad(q\in F).
\end{equation}

In the Final Remarks on page 45, Hatori asks whether a commutative
JB$^*$-triple always satisfies $(*)$.

\begin{figure}[ht]
  \centering
  \includegraphics[width=\textwidth]{figures/open_problem_crop.png}
  \caption{The open question in Hatori~\cite[p.~45]{Hatori}.}
  \label{fig:question}
\end{figure}

\section{Proof intuition}

An abelian JB$^*$-triple is a space $C^{\T}(L)$ of equivariant continuous
functions on a principal circle bundle.  Its relevant extreme functionals are
point evaluations, and choosing one representative per maximal-face class is
the same as choosing one point on every circle fiber.  Thus the representative
set in $(*)$ is a set-theoretic transversal of the bundle.

Condition $(*)$ prevents this transversal from acquiring a phase-shifted
limit.  Indeed, if representatives converge to $\lambda p$ with
$\lambda\ne1$, a relatively closed tail can be separated from $p$ by $(*)$;
weak$^*$ continuity then says its value at $\lambda p$ is small, contradicting
$p(a)=1$.  Hence the transversal is closed.  In a compact bundle it is then a
compact copy of the base, so it is a continuous global section.  Conversely,
a continuous section makes $(*)$ an immediate Urysohn-lemma construction.
The nontrivial Hopf bundle has no section and therefore supplies the desired
counterexample.

\section{Bundle characterization and counterexample}

Let $L$ be a compact principal $\T$-bundle, write $X=L/\T$, and put
\[
 E=C^{\T}(L):=\{a\in C(L):a(\lambda t)=\lambda a(t)
       \text{ for all }\lambda\in\T,\ t\in L\}.
\]
With its supremum norm and pointwise triple product, $E$ is a commutative
JB$^*$-triple.  For $t\in L$, write $\delta_t(a)=a(t)$.

We use the following standard face identification in precisely the form
recorded by Cabezas, Cueto-Avellaneda, Hirota, Miura, and
Peralta~\cite[Section~3]{CabezasEtAl}: the map
\begin{equation}\label{eq:eval}
 L\longrightarrow\cQ_E,\qquad t\longmapsto\delta_t,
\end{equation}
is a weak$^*$ homeomorphism in the compact case,
$\delta_{\lambda t}=\lambda\delta_t$, and the equivalence classes defining a
representative set are exactly the circle orbits.  Equivalently, a set of
representatives is a subset of $L$ meeting every fiber exactly once.

\begin{theorem}\label{thm:characterization}
For a compact principal circle bundle $\pi:L\to X$, the commutative
JB$^*$-triple $C^{\T}(L)$ satisfies Hatori's condition $(*)$ if and only if
$\pi$ admits a continuous global section.  Equivalently, $(*)$ holds if and
only if the principal bundle is trivial.
\end{theorem}

\begin{proof}
Suppose first that $(*)$ holds, witnessed by a representative set
$\cP\subset\cQ_E$.  We claim that $\cP$ is weak$^*$ closed in $\cQ_E$.
Let $r\in\overline{\cP}$.  By the representative property and
\eqref{eq:eval}, there are unique $p\in\cP$ and $\lambda\in\T$ such that
$r=\lambda p$.

Assume $\lambda\ne1$.  The compact Hausdorff space $\cQ_E$ has an open
neighborhood $U$ of $r$ whose closure does not contain $p$.  Set
\[
 F=\cP\cap\overline U.
\]
Then $F$ is relatively closed in $\cP$, $p\notin F$, and
$r\in\overline F$: every neighborhood of $r$ meets $\cP\cap U$ because
$r\in\overline{\cP}$.  Apply $(*)$ with $\varepsilon=1/2$.  There is
$a\in S(E)$ with
\[
 p(a)=1,\qquad |q(a)|\le\frac12\quad(q\in F).
\]
The function $q\mapsto q(a)$ is weak$^*$ continuous.  Since
$r\in\overline F$, it follows that $|r(a)|\le1/2$.  On the other hand,
$r=\lambda p$ gives $|r(a)|=|\lambda p(a)|=1$, a contradiction.  Hence
$\lambda=1$, so $r=p\in\cP$.  This proves that $\cP$ is closed.

Via \eqref{eq:eval}, regard $\cP$ as a closed transversal in $L$.  It is
compact and the restriction
\[
 \pi|_{\cP}:\cP\longrightarrow X
\]
is a continuous bijection.  Since $X$ is Hausdorff, this is a homeomorphism.
Its inverse is a continuous global section of $\pi$.

Conversely, let $s:X\to L$ be a continuous section.  It trivializes the
bundle by the homeomorphism
\[
 X\times\T\longrightarrow L,\qquad (x,\lambda)\longmapsto\lambda s(x).
\]
Choose the representatives $\cP=\{\delta_{s(x)}:x\in X\}$.  Let
$p=\delta_{s(x_0)}\in\cP$, and let $F\subset\cP$ be relatively closed with
$p\notin F$.  Under the homeomorphism $X\cong\cP$, the set $F$ corresponds
to a closed set $C\subset X$ not containing $x_0$.  By Urysohn's lemma there
is a continuous $h:X\to[0,1]$ such that $h(x_0)=1$ and $h|_C=0$.  Define
\[
 a(\lambda s(x))=\lambda h(x).
\]
Then $a\in S(C^{\T}(L))$, $p(a)=1$, and $q(a)=0$ for every $q\in F$.
Thus \eqref{eq:star} holds for every $\varepsilon>0$.
\end{proof}

\begin{theorem}\label{thm:hopf}
Let $\T$ act on
\[
 S^3=\{(z_1,z_2)\in\mathbb C^2:|z_1|^2+|z_2|^2=1\}
\]
by scalar multiplication.  The commutative JB$^*$-triple
\[
 C^{\T}(S^3)=
 \{a\in C(S^3):a(\lambda z)=\lambda a(z)\}
\]
does not satisfy Hatori's condition $(*)$.
\end{theorem}

\begin{proof}
The quotient map $S^3\to S^3/\T\cong S^2$ is the Hopf principal circle
bundle.  It has no continuous global section.  Indeed, a section would
trivialize the principal bundle and give a homeomorphism
$S^3\cong S^2\times S^1$, which is impossible because
\[
 \pi_1(S^3)=0,
 \qquad
 \pi_1(S^2\times S^1)\cong\mathbb Z
\]
(see, for example,~\cite{Hatcher}).  The conclusion follows from
Theorem~\ref{thm:characterization}.
\end{proof}

\section{Scope, verification, and novelty}

Theorem~\ref{thm:hopf} answers Hatori's question negatively.  The stronger
Theorem~\ref{thm:characterization} identifies the exact obstruction for all
compact principal-circle-bundle models: condition $(*)$ detects bundle
triviality.  This does not contradict the complex Mazur--Ulam theorem of
Cabezas et al.; in fact $C^{\T}(S^3)$ still has the complex Mazur--Ulam
property~\cite[Theorem~3.1]{CabezasEtAl}.  It shows that Hatori's sufficient
condition is not necessary, even within commutative JB$^*$-triples.

No computation is used.  The proof audit in \path{VERIFICATION.md} isolates
the face-representation input, the closed-transversal argument, and the Hopf
obstruction.

The local run indexes were searched for arXiv:2205.01888, the exact title,
Mazur--Ulam/Tingley terms, and condition $(*)$.  Bounded web searches on
August 9, 2026 used the exact question, title and author, ``commutative
JB$^*$-triple,'' ``condition $(*)$,'' principal circle bundles, and the Hopf
fibration.  They located the source and the known complex Mazur--Ulam theorem
but no later answer or this bundle characterization.  Novelty confidence is
positive but conservative pending expert literature review.

Human review is strongly recommended.  The decisive point is that $(*)$
forces an arbitrary representative transversal to be closed; once that is
confirmed, the compact-bundle topology makes the counterexample immediate.

\begin{thebibliography}{9}
\bibitem{Hatori}
O. Hatori,
\emph{The Mazur--Ulam property for a Banach space which satisfies a
separation condition}, RIMS K\^oky\^uroku Bessatsu \textbf{B93} (2023),
29--82; arXiv:2205.01888.

\bibitem{CabezasEtAl}
D. Cabezas, M. Cueto-Avellaneda, D. Hirota, T. Miura, and A. M. Peralta,
\emph{Every commutative JB$^*$-triple satisfies the complex Mazur--Ulam
property}, Ann. Funct. Anal. \textbf{13} (2022), Article 60;
arXiv:2201.06307; doi:10.1007/s43034-022-00204-6.

\bibitem{Hatcher}
A. Hatcher,
\emph{Algebraic Topology}, Cambridge University Press, 2002.
\end{thebibliography}

\end{document}
